% from ega0-1-fr.tex
\oldpage[0\textsubscript{I}]{26}

avec les conditions $\rho_U^U=\text{identité}$ et
$\rho_U^W=\rho_U^V\circ\rho_V^W$ si $U$, $V$, $W$ dans $\mathfrak{B}$ sont
tels que $U\subset V\subset W$. On peut lui associer un \emph{préfaisceau à
valeurs dans $\mathbf{K}$}~: $U\to\mathcal{F}'(U)$ au sens ordinaire, en
prenant pour tout ouvert $U$,
\[
  \mathcal{F}'(U)=\varprojlim_{V\subset U}\mathcal{F}(V),
\]
où $V$ parcourt l'ensemble ordonné (pour $\subset$, \emph{non filtrant en
général}) des ensembles $V\in\mathfrak{B}$ tels que $V\subset U$, car les
$\mathcal{F}(V)$ forment un système projectif pour les $\rho_V^W$
($V\subset W\subset U$, $V\in\mathfrak{B}$, $W\in\mathfrak{B}$). En effet,
si $U$, $U'$ sont deux ouverts de $X$ tels que $U\subset U'$, on définit
$\rho'_U{}^{U'}$ comme la limite projective (pour $V\subset U$) des morphismes
canoniques $\mathcal{F}'(U')\to\mathcal{F}(V)$, autrement dit l'unique
morphisme $\mathcal{F}'(U')\to\mathcal{F}'(U)$, qui, composé avec les
morphismes canoniques $\mathcal{F}'(U)\to\mathcal{F}(V)$, donne les morphismes
canoniques $\mathcal{F}'(U')\to\mathcal{F}(V)$~; la vérification de la
transitivité de $\rho'_U{}^{U'}$ est alors immédiate. De plus, si
$U\in\mathfrak{B}$, le morphisme canonique
$\mathcal{F}'(U)\to\mathcal{F}(U)$ est un isomorphisme, permettant d'identifier
ces deux objets%
\footnote{Si $X$ est un espace \emph{noethérien}, on peut encore définir
$\mathcal{F}'(U)$ et montrer que c'est un préfaisceau (au sens ordinaire)
lorsqu'on suppose seulement que $\mathbf{K}$ admet des limites projectives
pour les systèmes projectifs \emph{finis}. En effet, si $U$ est un ouvert
quelconque de $X$, il y a un recouvrement \emph{fini} $(V_i)$ de $U$ formé
d'ensembles de $\mathfrak{B}$~; pour tout couple $(i,j)$ d'indices, soit
$(V_{ijk})$ un recouvrement fini de $V_i\cap V_j$ formé d'ensembles de
$\mathfrak{B}$. Soit $I$ l'ensemble formé des $i$ et des triplets $(i,j,k)$,
ordonné par les seules relations $i>(i,j,k)$, $j>(i,j,k)$~; on prend alors
pour $\mathcal{F}'(U)$ la limite projective du système des $\mathcal{F}(V_i)$
et $\mathcal{F}(V_{ijk})$~; on vérifie aisément que cela ne dépend pas des
recouvrements $(V_i)$ et $(V_{ijk})$ et que $U\to\mathcal{F}'(U)$ est un
préfaisceau.}.
\end{env}

\begin{env}[3.2.2]
\label{0.3.2.2-fr}
Pour que le préfaisceau $\mathcal{F}'$ ainsi défini soit un \emph{faisceau},
il faut et il suffit que le préfaisceau $\mathcal{F}$ sur $\mathfrak{B}$
vérifie la condition~:

\smallskip
\noindent
\textup{(F\textsubscript{0})} \emph{Pour tout recouvrement $(U_\alpha)$ de
$U\in\mathfrak{B}$ par des ensembles $U_\alpha\in\mathfrak{B}$ contenus dans
$U$, et pour tout objet $T\in\mathbf{K}$, l'application qui, à tout
$f\in\operatorname{Hom}(T,\mathcal{F}(U))$ fait correspondre la famille
\[
  (\rho_{U_\alpha}^{U}\circ f)
  \in\prod_\alpha\operatorname{Hom}(T,\mathcal{F}(U_\alpha))
\]
est une bijection de $\operatorname{Hom}(T,\mathcal{F}(U))$ sur l'ensemble des
$(f_\alpha)$ tels que
$\rho_V^{U_\alpha}\circ f_\alpha=\rho_V^{U_\beta}\circ f_\beta$ pour tout
couple d'indices $(\alpha,\beta)$ et tout $V\in\mathfrak{B}$ tel que
$V\subset U_\alpha\cap U_\beta$}%
\footnote{Cela signifie encore que le couple formé par $\mathcal{F}(U)$ et les
$\rho_\alpha=\rho_{U_\alpha}^{U}$ est \emph{solution du problème universel}
défini (3.1.1) par la donnée de
$A_\alpha=\mathcal{F}(U_\alpha)$,
$A_{\alpha\beta}=\prod\mathcal{F}(V)$ (pour les $V\in\mathfrak{B}$ tels que
$V\subset U_\alpha\cap U_\beta$) et
$\rho_{\alpha\beta}=(\rho_V''):\mathcal{F}(U_\alpha)\to
\prod\mathcal{F}(V)$ défini par la condition que pour
$V\in\mathfrak{B}$, $V'\in\mathfrak{B}$, $W\in\mathfrak{B}$,
$V\cup V'\subset U_\alpha\cap U_\beta$, $W\subset V\cap V'$,
$\rho_W^V\circ\rho_V''=\rho_W^{V'}\circ\rho_{V'}''$.}.

La condition est évidemment nécessaire. Pour montrer qu'elle est suffisante,
considérons d'abord une seconde base $\mathfrak{B}'$ de la topologie de $X$,
\emph{contenue dans $\mathfrak{B}$}, et montrons que si $\mathcal{F}''$
désigne le préfaisceau déduit de la sous-famille
$(\mathcal{F}(V))_{V\in\mathfrak{B}'}$, $\mathcal{F}''$ est
\emph{canoniquement isomorphe} à $\mathcal{F}'$. En effet, tout d'abord la
limite projective (pour $V\in\mathfrak{B}'$, $V\subset U$) des morphismes
canoniques $\mathcal{F}'(U)\to\mathcal{F}(V)$ est un morphisme
$\mathcal{F}'(U)\to\mathcal{F}''(U)$ pour tout ouvert $U$. Si
$U\in\mathfrak{B}$, ce morphisme est un isomorphisme, car par hypothèse les
morphismes canoniques $\mathcal{F}''(U)\to\mathcal{F}(V)$ pour
$V\in\mathfrak{B}'$, $V\subset U$, se factorisent en
\[
  \mathcal{F}''(U)\to\mathcal{F}(U)\to\mathcal{F}(V),
\]
et il est immédiat de voir que les composés des morphismes
$\mathcal{F}(U)\to\mathcal{F}''(U)$ et
$\mathcal{F}''(U)\to\mathcal{F}(U)$ ainsi définis sont les identités. Ceci
étant, pour tout ouvert $U$, les morphismes
$\mathcal{F}''(U)\to\mathcal{F}''(W)=\mathcal{F}(W)$ pour
$W\in\mathfrak{B}$ et $W\subset U$ vérifient les conditions caractérisant la
limite projective des $\mathcal{F}(W)$ ($W\in\mathfrak{B}$, $W\subset U$), ce
qui démontre notre assertion compte tenu de l'unicité d'une limite projective
à un isomorphisme près.

Cela posé, soit $U$ un ouvert quelconque de $X$, $(U_\alpha)$ un recouvrement
de $U$ par des ouverts contenus dans $U$, et soit $\mathfrak{B}'$ la
sous-famille de $\mathfrak{B}$ constituée par les ensembles

